\subsection{Symplectic group.}

Suppose the ring A is commutative. If \(\Phi\) is an alternating bilinear form on E, the automorphisms of the module E leaving \(\Phi\) invariant are called symplectic automorphisms (or symplectic transformations) relative to \(\Phi\) , and they form a group called the symplectic group associated with \(\Phi\) ; it is sometimes denoted \(\mathbf{Sp}(\Phi)\) .

Consider in particular, on the module \(E = A^{2m}\) the alternating bilinear form \(\Phi_{0}\) whose matrix with respect to the canonical basis \((e_{i})\) of E is

\[
R _ {m} = \left( \begin{array}{c c} 0 & I _ {m} \\ - I _ {m} & 0 \end{array} \right).
\]

The symplectic automorphisms and the symplectic group relative to \(\Phi_{0}\) are simply called symplectic automorphisms and symplectic group in \(2m\) variables (over A), respectively; this group is denoted \(\mathbf{Sp}(2m, A)\) or \(\mathbf{Sp}_{2m}(A)\) . Every matrix \(A\) of a symplectic automorphism with respect to the canonical basis \((e)_i\) is called a symplectic matrix. Such a matrix is invertible, and, by formula (48) of § 1, no. 10, satisfies the relation

\[
{ } ^ { t } A . R _ { m } . A = R _ { m } .\tag{9}
\]

Conversely, if a square matrix A of order 2m over A satisfies (9), it is symplectic: it suffices indeed to prove that it is invertible; but (9) implies Pf(R \(_{m}\) ) = Pf( \(^{t}\) A . R \(_{m}\) . A) = (det A)Pf(R \(_{m}\) ) by prop. 1 of no. 2, hence det A = 1. We have at the same time proved the following proposition:

PROPOSITION 3. --- The determinant of a symplectic matrix is equal to 1.

If A is a commutative field and \(\Phi\) a nondegenerate alternating bilinear form on a vector space E of even dimension \(2m\) over A, the symplectic group associated with \(\Phi\) is isomorphic to \(\mathbf{Sp}(2m, \mathrm{A})\) , by cor. 1 of th. 1.

Exercises. --- ¶ 1) Let A be a principal commutative ring, E a free A-module of finite dimension n, Φ an alternating bilinear form on E; the ideals Aαi (1 ≤ i ≤ r) defined in th. 1 of no. 1 are called the invariant factors of Φ.

\begin{enumerate}
\def\labelenumi{\alph{enumi})}
\item
  Let F be a submodule of E, and let \(\Phi_{\mathrm{F}}\) the restriction of \(\Phi\) to F × F. Show that if A \(\beta_{i}\) ( \(1 \leqslant i \leqslant s\) ) are the invariant factors of \(\Phi_{\mathrm{F}}\) (where \(\beta_{i}\) divides \(\beta_{i+1}\) ), we have \(s \leqslant r\) and \(\beta_{i}\) is a multiple of \(\alpha_{i}\) for \(1 \leqslant i \leqslant s\) . (Reduce to the case where \(r = s = n/2\) , and use exerc. 9 b) and 9 c) of chap. VII, § 4).
\item
  Let \(E_{1}\) be a second free A-module of finite dimension, \(\Phi_{1}\) an alternating bilinear form on \(E_{1}, A\gamma_{1}, \ldots, A\gamma_{s}\) its invariant factors ( \(\gamma_{i}\) dividing \(\gamma_{i+1}\) ). In order that \(\Phi_{1}\) be the inverse image of \(\Phi\) by a linear mapping of \(E_{1}\) into E, it is necessary and sufficient that \(s \leqslant r\) and \(\gamma_{i}\) be a multiple of \(\alpha_{i}\) for \(1 \leqslant i \leqslant s\) . (Use a) and prop. 4 of chap. VII, § 4, no. 5).
\item
  Let F, G be two submodules of E, such that \(F^{0}\) (resp. \(G^{0}\) ) be a supplement of F (resp. G) in E. If the restrictions of \(\Phi\) to F and G are equivalent, show that the same holds for the restrictions of \(\Phi\) to \(F^{0}\) and to \(G^{0}\) , and that there exists an automorphism of E leaving \(\Phi\) invariant and mapping F to G.
\item
  Give an example of two submodules F, G of E, of dimension 2, such that F and G admit supplements in E and such that the restrictions \(\Phi_{\mathrm{F}}\) and \(\Phi_{\mathrm{G}}\) are equivalent, but that there exists no automorphism of E leaving \(\Phi\) invariant and mapping F to G (take \(n = 4\) ).
\end{enumerate}

\begin{enumerate}
\def\labelenumi{\arabic{enumi})}
\setcounter{enumi}{1}
\tightlist
\item
  Let \(\Phi\) be an alternating bilinear form on a finite-dimensional vector space E. Show that for every subspace M of E, the difference dim M - dim (M ∩ M⁰) is even. (First consider the case where Φ is nondegenerate).
\end{enumerate}

¶ 3) Let E be a vector space of even dimension n = 2m over a commutative field A; let Φ and Ψ be two alternating bilinear forms on E; suppose Ψ is nondegenerate. Let u and v be the linear maps from E to E* associated on the right with Φ and Ψ respectively; v is an isomorphism from E onto E*; set w = v⁻¹ ∘ u, so that w is an endomorphism of E.

\begin{enumerate}
\def\labelenumi{\alph{enumi})}
\item
  Set \(\mathbf{M}_{\mathbf{0}} = \mathbf{E}\) and by induction \(\mathbf{M}_{k+1} = w(\mathbf{M}_{k})\) for \(k > 0\) . Show that if \(\mathbf{M}_{k}^{\prime}\) is the orthogonal subspace to \(\mathbf{M}_{k}\) for \(\Phi\) , \(\mathbf{M}_{k+1}\) is the orthogonal subspace to \(\mathbf{M}_{k}^{\prime}\) for \(\Psi\) .
\item
  Let \(n_{0}\) be the dimension of \(\mathbf{M}_{0}^{\prime}\) ; we set \(m_{0}=0\) , and for \(k\geqslant1\) , we denote by \(m_{k}\) the dimension of \(\mathbf{M}_{k}\cap\mathbf{M}_{0}^{\prime}\) . Show that for \(k\geqslant1\) the dimension of \(\mathbf{M}_{k}\) is \(n-n_{0}-(m_{1}+\ldots+m_{k-1})\) and that of \(\mathbf{M}_{k}^{\prime}\) is \(n_{0}+(m_{1}+\ldots+m_{k})\) .
\item
  Show that, for all \(k \geqslant 0\) , the dimension of \(\mathbf{M}_k \cap \mathbf{M}_k'\) is \(m_k + m_{k+1} + \ldots + m_{2k}\) , and that of \(\mathbf{M}_k' \cap \mathbf{M}_{k+1}\) is \(m_{k+1} + m_{k+2} + \ldots + m_{2k+1}\) (apply exerc. 2 b) of § 3 to compute the dimensions of \(\mathbf{M}_h \cap \mathbf{M}_{2k-h}'\) and of \(\mathbf{M}_h' \cap \mathbf{M}_{2k+1-h}\) by induction on \(h\) ).
\item
  Deduce from c) that the numbers \(m_{k}\) are even (use exerc. 2).
\item
  Conclude from d) that the number of elementary divisors of the endomorphism w, corresponding to the characteristic root \(\lambda = 0\) , and having a given degree, is even (cf. chap. VII, § 5, exerc. 20).
\end{enumerate}

\begin{enumerate}
\def\labelenumi{\arabic{enumi})}
\setcounter{enumi}{3}
\item
  Let E be a vector space over a commutative field A, admitting a countable basis \((e_n)_{n\geqslant 1}\) , \(\Phi\) a nondegenerate alternating form on E. Show that there exists in E a basis \((a_n)\) such that \(\Phi(a_{2n-1}, a_{2n}) = 1\) for all \(n\geqslant 1\) , and \(\Phi(a_i, a_j) = 0\) for any other pair of indices such that \(i < j\) (reason as in exercise 13 of § 4).
\item
  For any alternating matrix \(X = (x_{ij})\) of even order \(n = 2m\) over a commutative ring, and for every index \(i\) , show that one has \(\operatorname{Pf}(X) = \sum_{j=1}^{n} (-1)^{i+j-1} \operatorname{Pf}(X_{ij}) x_{ij}\) , where \(X_{ij}\) is the matrix of order \(n - 2\) obtained by deleting from \(X\) the rows and columns with indices \(i\) and \(j\) .
\item
  Let M be a square matrix of order m over a commutative ring. Show that if we set
\end{enumerate}

\[
R = \left( \begin{array}{c c} 0 & M \\ - ^ {t} M & 0 \end{array} \right)
\]

we have \(\operatorname{Pf}(R) = \det M\) (first prove it when \(M\) is invertible, using formula (7)).

\begin{enumerate}
\def\labelenumi{\arabic{enumi})}
\setcounter{enumi}{6}
\item
  Let A be a commutative field, \(\mathfrak{A}(A)\) the set of alternating matrices of order \(2m\) on A; let I be a mapping from \(\mathfrak{A}(A)\) into A such that for every matrix \(R\in\mathfrak{A}(A)\) and for every matrix \(P\) of order \(2m\) over A, one has \(\mathrm{I}(^{t}PRP)=(\det P)^{h}\mathrm{I}(R)\) , where \(h\) is a rational integer. Show that \(\mathrm{I}(R)=c(\mathrm{Pf}R)^{h}\) , where \(c\in\mathrm{A}\) (using formula (7) and th. 1 of no. 1).
\item
  Let \(P\) , \(Q\) be two alternating square matrices of even order \(2m\) on a commutative ring A. Let \(\varphi(X) = \mathrm{Pf}(P - XQ)\) ; show that, if \(Q\) is invertible, one has \(\varphi(Q^{-1}P)=0\) . (First consider the case where A is a field; then deduce from exerc. 3 that the minimal polynomial of the matrix \(Q^{-1}P\) divides \(\varphi(X)\) , by passing to an algebraically closed extension of A, and noting that \(\varphi^{2}\) is the characteristic polynomial of \(Q^{-1}P\) up to a scalar factor).
\end{enumerate}

¶ 9) Let E be a vector space of dimension n over a commutative field A, Φ an alternating bilinear form on E, Ψ a sesquilinear Hermitian form (for an involutive automorphism of A) on E, and P and Q the respective matrices of Φ and Ψ with respect to the same basis of E.

\begin{enumerate}
\def\labelenumi{\alph{enumi})}
\item
  We assume \(\Psi\) nondegenerate. Show that the number of elementary divisors of \(Q^{-1}P\) corresponding to the characteristic root 0 and having the same even degree, is an even number (method of exerc. 3).
\item
  We assume that \(\Phi\) be nondegenerate (which implies that \(n\) is even). Show that the number of elementary divisors of \(P^{-1}Q\) corresponding to the characteristic root 0, and having the same odd degree, is an even number (same method).
\end{enumerate}

\begin{enumerate}
\def\labelenumi{\arabic{enumi})}
\setcounter{enumi}{9}
\tightlist
\item
  Let \(\omega_{2m}(\mathbf{F}_q)\) be the order of the symplectic group \(\mathbf{Sp}(2m, \mathbf{F}_q)\) over the finite field \(\mathbf{F}_q\) . Show that if \(h_{2m}\) is the number of pairs of vectors \((x, y)\) of \(\mathbf{F}_q^{2m}\) such that \(\Phi_0(x, y) = 1\) (notations of no. 3), one has \(\omega_{2m}(\mathbf{F}_q) = h_{2m}\omega_{2m-2}(\mathbf{F}_q)\) ; deduce from this
\end{enumerate}

\[
\omega_ {2 m} (\mathbf {F} _ {q}) = (q ^ {2 m} - 1) q ^ {2 m - 1} (q ^ {2 m - 2} - 1) q ^ {2 m - 3} \dots (q ^ {2} - 1) q.
\]

\begin{enumerate}
\def\labelenumi{\arabic{enumi})}
\setcounter{enumi}{10}
\tightlist
\item
  Let A be a commutative field. Show that every transformation \(u\) belonging to the symplectic group \(\mathbf{Sp}(2m, \mathrm{A})\) is a product of transvections belonging to this group (called symplectic transvections; cf.§ 4, exerc. 6). (Reason by induction on \(m\) , by showing that if \(x, y\) are two non-orthogonal vectors of \(\mathrm{E} = \mathrm{A}^{2m}\) , there exists a product \(\nu\) of symplectic transvections such that \(\nu u\) leaves \(x\) and \(y\) invariant). Deduce a new proof of prop. 3 of no. 3.
\end{enumerate}

*12) Let A be a commutative field, E a vector space of even dimension \(n=2m\) over A, \(\Phi\) a nondegenerate alternating bilinear form on E. Show that, for every similarity \(u\) for the form \(\Phi\) ( \(\S 6\) , no. 5), with multiplier \(\alpha\) , we have det \(u=\alpha^{m}\) (using formula (7)).*

¶ 13) We assume that A is a commutative field of characteristic 0, E a vector space of dimension 2m over A, Φ a nondegenerate alternating bilinear form on E. We identify the inverse form \(\widehat{\Phi}\) of Φ with a bivector \(\Gamma\in\bigwedge E\) , so that for any symplectic basis \((e_{i})_{1\leqslant i\leqslant2m}\) of E (for \(\Phi\) ), indexed such that \(\Phi(e_{i},e_{j})=\Phi(e_{m+i},e_{m+j})=0\) , \(\Phi(e_{i},e_{m+j})=\delta_{ij}(1\leqslant i\leqslant m,1\leqslant j\leqslant m)\) , one has

\[
\Gamma = e _ {1} \wedge e _ {m + 1} + e _ {2} \wedge e _ {m + 2} + \dots + e _ {m} \wedge e _ {2 m}.
\]

We say that a nonzero decomposable \(p\) -vector \(z \in \bigwedge E\) is isotropic (resp. totally isotropic) for \(\Phi\) if the vector subspace \(V_z\) corresponding to \(z\) (chap. III, § 7, no. 3) is isotropic (resp. totally isotropic).

\begin{enumerate}
\def\labelenumi{\alph{enumi})}
\item
  If z is a decomposable p-vector \(\neq0\) , 2r the dimension of a supplement of \(V_{z}\cap V_{2}^{0}\) with respect to \(V_{z}\) , show that \(m-p+r\) is the largest of the integers h such that one has \(z\wedge\Gamma^{h}\neq0\) , denoting by \(\Gamma^{h}\) the h-th power of \(\Gamma\) in the exterior algebra \(\wedge E\) (using prop. 2 of § 4, no. 2).
\item
  Show that every bivector \(x \in \bigwedge E\) can be written in the form \(\lambda \Gamma + x_{1}\) , where \(\lambda\) is a scalar and \(x_{1}\) a linear combination of totally isotropic decomposable bivectors. (Reduce to the case where \(x\) is decomposable, and note that if \((e_{i})\) is a symplectic basis, \((e_{1} + e_{2}) \wedge (e_{m+1} - e_{m+2})\) is totally isotropic).
\item
  If \(p \leqslant m\) , show that every \(p\) -vector \(z \in \hat{\Lambda}\) E can be written \(z = x \wedge \Gamma + z_1\) , where \(x\) is a \((p - 2)\) -vector and \(z_1\) a linear combination of \(p\) -vectors that are decomposable and totally isotropic. (Reduce to the case where \(z\) is decomposable, and reason by induction on \(p\) . One can thus reduce to the case where, \((e_i)\) being a symplectic basis, one has \(z = e_1 \wedge e_2 \wedge \ldots \wedge e_{p-1} \wedge e_{m+p-1}\) , and consider the bivectors \((e_{p-1} + e_{p+i}) \wedge (e_{m+p-1} - e_{m+p+i})\) for \(0 \leqslant i \leqslant m - p\) ).
\item
  For \(1 \leqslant p \leqslant m\) , show that every \((m+p)\) -vector \(z \in \wedge\) E can be written \(z = y \wedge \Gamma^p\) , where \(y\) is a \((m-p)\) -vector. (Reduce to the case where \(z\) is decomposable; if \(2r\) is the dimension of a supplement of \(\mathrm{V}_z \cap \mathrm{V}_z^0\) with respect to \(\mathrm{V}_z\) , distinguish two cases, according as \(r = p\) or \(r > p\) ; in the second case, reason by induction on \(r\) , in the same manner as in \(c)\) . Deduce from this that the map \(y \to y \wedge \Gamma^p\) of \(\wedge^{m-p}\) E into \(\wedge^{m+p}\) E is bijective (note that the two spaces have the same dimension). If y is a decomposable \((m-p)\) -vector, show that \(z = y \wedge \Gamma^{p}\) is decomposable and that \(V_{z} = V_{y}^{0}\) .
\end{enumerate}

\begin{enumerate}
\def\labelenumi{\alph{enumi})}
\setcounter{enumi}{4}
\tightlist
\item
  Deduce from d) that, for \(p \leqslant m\) , the subspace \(\hat{E}\) is the direct sum of the homogeneous component of degree p of the two-sided ideal c generated by \(\Gamma\) in \(\hat{E}\) , and of the subspace \(R_{p}\) of p-vectors \(z_{1}\) such that
\end{enumerate}

\(z_{1} \wedge \Gamma^{m-p+1} = 0\) (for \(z \in \bigwedge E\) , apply \(d\) ) to the \((2m - p + 2)\) -vector \(z \wedge \Gamma^{m-p+1}\) ). Using \(c\) ), prove that \(R_{p}\) is generated by the decomposable totally isotropic \(p\) -vectors, and by using \(d\) ), show

that the map \(x \to x \wedge \Gamma\) of \(\wedge E\) into \(\wedge E\) is injective, and that \(R_p\) is of dimension \(\binom{2m}{p} - \binom{2m}{p-2}\) .

\begin{enumerate}
\def\labelenumi{\alph{enumi})}
\setcounter{enumi}{5}
\tightlist
\item
  Show that if \(p \leqslant m\) , a necessary and sufficient condition for a p-vector z to be of the form \(x \wedge \Gamma\) is that, for every decomposable totally isotropic m-vector u, one has \(z \wedge u = 0\) . (Show that, if this condition is satisfied, and if \(z \in R_{p}\) , one has \(z \wedge y = 0\) for all \((2m - p)\) -vector y; for this, put y in the form \(x \wedge \Gamma^{m-p}\) where x is a p-vector, then apply c) to x, and note that if \(x_{1}\) is a decomposable totally isotropic p-vector, \(x_{1} \wedge \Gamma^{m-p}\) is a decomposable \((2m - p)\) -vector which can be written \(u_{1} \wedge \varphi_{1}\) , where \(u_{1}\) is a decomposable totally isotropic \(m\) -vector).
\end{enumerate}

Let a be the annihilator of the ideal c in \(\wedge\) E; a is the direct sum of the subspaces \(R_{m}\) , \(R_{m-1} \wedge \Gamma, \ldots, R_{1} \wedge \Gamma^{m-1}\) and \(K\Gamma^{m}\) . Show that the annihilator of a in \(\wedge\) E is equal to c (cf.§ 2, exerc. 4 b)).

\begin{enumerate}
\def\labelenumi{\alph{enumi})}
\setcounter{enumi}{6}
\tightlist
\item
  For every symplectic automorphism u of E (for \(\Phi\) ), let \(\bar{u}\) be the canonical extension of u to an automorphism of the algebra \(\wedge\) E (chap. III, § 5, no. 9). Show that the only elements of \(\wedge\) E invariant under all the automorphisms \(\bar{u}\) are the linear combinations of 1, \(\Gamma\) , \(\Gamma^{2}\) , \(\ldots\) , \(\Gamma^{m}\) with coefficients in K. (If ( \(e_{i}\) ) is a symplectic basis, write that an element of \(\wedge\) E is invariant under the symplectic transvections (exerc. 11) corresponding to the hyperplanes orthogonal to the \(e_{i}\) ; then consider the symplectic transformations \(u_{ij}\) defined by \(u_{ij}(e_{i}) = e_{j}\) , \(u_{ij}(e_{j}) = e_{i}\) , \(u_{ij}(e_{m+i}) = e_{m+j}\) , \(u_{ij}(e_{m+j}) = e_{m+i}\) , \(u_{ij}(e_{k}) = e_{k}\) for every other index k).
\end{enumerate}

\begin{enumerate}
\def\labelenumi{\arabic{enumi})}
\setcounter{enumi}{13}
\item
  Let A be a commutative field of characteristic 2, E the vector space \(\mathrm{A}^{2m}\) ; we identify the symplectic group \(\mathbf{Sp}(2m,\mathrm{A})\) with the group of symplectic matrices \(U\) , thus satisfying the relation \(^t U.R.U = R\) , where \(R = \left( \begin{array}{cc}0 & I_m\\ I_m & 0 \end{array} \right)\) . We set \(D = \left( \begin{array}{cc}0 & 0\\ I_m & 0 \end{array} \right)\) ; show that every symplectic matrix \(U\) such that \(I + U\) is invertible can be written in a unique way in the form \((^t D + S)^{-1}(D + S)\) , where \(S\) is a symmetric matrix such that \(D + S\) is invertible, and conversely (cf.§ 4, exerc. 11).
\item
  Let A be a commutative field, E a vector space of dimension 2m over A, Φ a nondegenerate alternating bilinear form on E. a) Extend to endomorphisms u of E such that \(u^{*}u = uu^{*}\) ( \(u^{*}\) being the adjoint of u with respect to Φ) the results of § 4, exerc. 12 a) and b).
\end{enumerate}

\begin{enumerate}
\def\labelenumi{\alph{enumi})}
\setcounter{enumi}{1}
\item
  We assume that \(u^{*} = u\) . With the notations of exerc. 12 of § 4, show that if M is a minimal element of the set \(\mathfrak{M}\) of non-isotropic subspaces contained in \(G(p, p)\) and stable under \(u\) , then either M is an indecomposable submodule of \(E_{u}\) , or it is a direct sum of two such isomorphic submodules (reason as in exerc. 12 c) of § 4). Show by examples (with \(p(X) = X - 1\) ) that the two cases can occur.
\item
  We assume \(u^* = -u\) , with A not of characteristic 2. With the same notations, let \(p^h\) be the minimal polynomial of the restriction of \(u\) to M, and let \(d\) be the degree of \(p\) . Show that if \(d(h - 1)\) is odd, M is an indecomposable submodule of \(E_u\) ; if on the contrary \(d(h - 1)\) is even, M is either indecomposable or a direct sum of two isomorphic indecomposable submodules.
\item
  We assume that \(u^{*}u = 1\) (in other words that \(u \in \mathbf{Sp}(\Phi)\) ). With the notations of c), if \(p(X)\) does not divide \(X^{d(h-1)} - 1\) , or if \(p(X) = X - 1\) and \((-1)^{h} \neq -1\) , M is an indecomposable submodule of \(E_{u}\) ; otherwise, M is either indecomposable or a direct sum of two isomorphic indecomposable submodules.
\end{enumerate}

